
\question
One type of webbed toes is a Y-linked trait. (Y-linked traits are found only on the Y chromosome.) If a male with webbed toes marries a female that does not have the trait, their children might have which of the following phenotypes?

\begin{choices}
\choice Normal females, web-toed females, normal males, and web-toed males.
\choice Normal females, normal males or web-toed males.
\choice Normal females, web-toed females, and web-toed males.
\choice Normal females, web-toed females, and normal males.
\correctchoice Normal females and web-toed males.
\end{choices}

\question
Fragile X syndrome is a sex-linked dominant trait that causes a range of developmental problems including learning disabilities and cognitive impairment.  Usually, the disease is worse in males than females.  Mary has the trait but her husband Mark does not. Their daughter, Michelle, does not have the trait. What is the probability that their next two children will both have the trait?

\begin{choices}
\choice 75\%
\choice 50\%
\choice 100\%
\correctchoice 25\%
\choice 12.5\%
\end{choices}

\question
After using the polymerase chain reaction to amplify the amelogenin gene from a blood sample, you find two distinct peaks (or bands) in the gel.  What can you conclude?

\begin{choices}
\correctchoice The blood sample came from a male.
\choice The blood sample came from the victim.
\choice The blood sample came from a female.
\choice The blood sample did not come from a human.
\choice The blood sample came from the suspect.
\end{choices}

\question
The sugary gene contributes to the sweetness of corn. A corn plant that is \textit{Ss} for the sugary gene is crossed with a plant that is \textit{ss}. What is the expected genotypic ratio of offspring?

\begin{choices}
\correctchoice 1 \textit{Ss} : 1 \textit{ss}
\choice 1 \textit{SS} : 1 \textit{Ss} : 1 \textit{ss}
\choice 1 \textit{SS} : 1 \textit{ss}
\choice 1 \textit{SS} : 2 \textit{Ss} : 1 \textit{ss}
\choice 1 \textit{SS} : 1 \textit{Ss}
\end{choices}

\newpage

\question
Danny Ramos Gomez, his brother Larry, and his sister Jamie, all have hypertrichosis. Their disease is a sex-linked, dominant trait with alleles \textit{H} and \textit{h}. Based on his family pedigree, which of the following will be true for Danny's daughter and Larry's son?

\begin{multicols}{2}
\begin{choices}
\choice Danny's daughter and Larry's son both will have hypertrichosis.
\choice Danny's daughter and Larry's son both will not have hypertrichosis.
\correctchoice Danny's daughter will have hypertrichosis. Larry's son will not have hypertrichosis.
\choice Danny's daughter may or many not have hypertrichosis, depending on the genotype of his wife. Larry's son will have hypertrichosis.
\choice Danny's daughter may or many not have hypertrichosis, depending on the genotype of his wife. Larry's son will not have hypertrichosis.
\end{choices}
\columnbreak
\includegraphics[width=0.4\textwidth]{gomez_pedigree}
\end{multicols}

\question
Assume you cross a black-feathered chicken with a white-feathered chicken, then cross two offspring from the F$_1$ generation.  In the F$_2$ generation, $1/4$ of the chickens have black feathers, $1/4$ of the chickens have white feathers, and $1/2$ of the chickens are speckled.  When you look closely, you notice that the speckles are caused by equal numbers of white and black feathers on the body. This type of inheritance is called

\begin{choices}
\correctchoice codominance.
\choice environmental effects.
\choice incomplete dominance.
\choice continuous dominance.
\choice complete dominance.
\end{choices}

\question
The two strands of DNA are held together by

\begin{choices}
\choice three bonds between A and G, and two bonds between C and T.
\choice three bonds between C and G, and three bonds between A and T.
\choice three bonds between A and T, and two bonds between C and G.
\correctchoice three bonds between C and G, and two bonds between A and T.
\choice two bonds between C and T, and two bonds between A and G.
\end{choices}

\question
The generation that results from crosses between true-breeding parents is called

\begin{choices}
\choice the P$_1$ generation
\choice the true-breeding generation.
\correctchoice the F$_1$ generation.
\choice the inbred generation.
\choice the F$_2$ generation.
\end{choices}

\question
Compare the DNA sequence of the parent cell to the sequence of a daughter cell.  Which one type of mutation best explains this difference?

\vspace{0.4\baselineskip}
\begin{tabular}{@{}ll@{}}
Parent: & $5'$ -- \texttt{CCGTTGAAGTACATTGAAGCTTT} -- $3'$\\
Daughter: & $5'$ -- \texttt{CCGTTGAAGTATATTGAAGCTTT} -- $3'$\\
\end{tabular}

\begin{choices}
\correctchoice point mutation
\choice gene duplication
\choice jumping gene
\choice transposable element
\choice frameshift mutation
\end{choices}

\question
After a plant breeder crossed two pink carnations, 24\% of the offspring had red flowers, 54\% had pink flowers, and 22\% had white flowers. According to these results, both alleles are most likely

\begin{choices}
\correctchoice incompletely dominant.
\choice dominant.
\choice recessive.
\choice sex-linked.
\choice codominant.
\end{choices}

\question
Which sequence below would make a good PCR primer for the following target gene?

\vspace{0.4\baselineskip}
%\begin{tabular}{@{}ll@{}}
Target gene: $3'$ -- \texttt{TCGGAATCTCTGGTCAATGCC} -- $5'$
%\vspace{0.4\baselineskip}

\begin{choices}
\choice $5'$ -- \texttt{GATTCCG} -- $3'$
\choice $3'$ -- \texttt{TCGGAAT} -- $5'$
\correctchoice $5'$ -- \texttt{AGCCTTA} -- $3'$
\choice $3'$ -- \texttt{AGCCTTA} -- $5'$
\choice $3'$ -- \texttt{GATTCCG} -- $5'$
\choice $5'$ -- \texttt{TCGGAAT} -- $3'$
\end{choices}

\question
Which of the following is not a nucleotide?

\begin{choices}
\choice Guanine
\correctchoice Urocine
\choice Adenine
\choice Thymine
\choice Cytosine
\end{choices}

\question
The \underline{\hspace{2cm}} strands of DNA form a \underline{\hspace{2cm}} held together by \underline{\hspace{2cm}} bonds.

\begin{choices}
\correctchoice two, double helix, hydrogen.
\choice one, single helix, hydrogen.
\choice two, double helix, covalent.
\choice one, double helix, nitrogen.
\choice two, single helix, carbon.
\end{choices}

\question
A change in the nucleotide sequence of one or more genes is called

\begin{choices}
\correctchoice mutation.
\choice recombination.
\choice phenotype.
\choice deletion (like the amelogenin gene).
\choice genotype.
\end{choices}

\question
The kernals on ears of corn come in a variety of colors and textures. The \textit{P} gene controls color.  The dominant \textit{P} allele produces purple kernals. The recessive \textit{p} allele produces yellow kernals.  The \textit{S} gene controls texture. The dominant \textit{S} allele produces smooth kernals. The recessive \textit{s} allele produces wrinkled kernals. If you cross a \textit{ppss} individual with a \textit{PpSs} individual, the expected genotype ratio for corn kernals in the next generation is 

\begin{choices}
\choice  50\% \textit{PpSs}: 25\% \textit{Ppss}: 25\% \textit{ppss}
\choice 25\% \textit{ppSs}: 50\% \textit{Ppss}: 25\% \textit{ppss}
\correctchoice 25\% \textit{PpSs}: 25\% \textit{Ppss}: 25\% \textit{ppSs}: 25\% \textit{ppss}
\choice  25\% \textit{PpSs}: 50\% \textit{ppSs}: 25\% \textit{Ppss}
\choice 9 \textit{PpSs} : 3 \textit{Ppss} : 3 \textit{ppSs} : 1 \textit{ppss}
\end{choices}

\question
The one allele that determines the phenotype of a heterozygote is

\begin{choices}
\choice the recessive gene.
\correctchoice the dominant allele.
\choice the allele on the dominant gene.
\choice the dominant gene.
\choice the recessive allele.
\end{choices}

\question
DNA is a long, chainlike molecule made up of subunits called

\begin{choices}
\choice genes.
\choice proteins.
\choice base pairs.
\correctchoice nucleotides.
\choice helices (plural of helix).
\end{choices}

\question
In a cross between two heterozygotes, the Punnett square shows the three possible \underline{\hspace{2cm}} and two possible \underline{\hspace{2cm}} that may appear in the offspring.

\begin{choices}
\choice allele; genotypes
\correctchoice genotypes; alleles
\choice phenotypes; genes
\choice phenotypes; genotypes
\choice genotypes; phenotypes
\end{choices}

\newpage

\question
Dogs have a hereditary form of deafness caused by the recessive \textit{d} allele (not sex-linked).  Betty has a dog that can hear but she does not know whether the dog is \textit{DD} or \textit{Dd}.  She breeds her dog to one that is deaf.  If Betty's dog is homozygous, what is the probability of producing deaf puppies?

\begin{choices}
\choice 75\%.
\choice 100\%.
\choice 25\%.
\choice 50\%.
\correctchoice 0\%.
\end{choices}

\question
Mutations which occur in 

\begin{choices}
\choice germ-line and somatic tissues cannot be passed on to future generations.
\choice germ-line  and somatic tissues can be passed on to future generations.
\correctchoice germ-line tissues can be passed on to future generations.
\choice germ-line tissues cannot be passed on to future generations.
\choice somatic tissues can passed on to future generations.
\end{choices}

\question
In cats, the agouti gene (\textit{A}/\textit{a}) determines the color of each hair. The dominant \textit{A} allele produces orange hairs with black tips.  The recessive \textit{a} allele produces black hairs.  The tabby gene (\textit{M}/\textit{m}) produces one of two tabby patterns.  The dominant \textit{M} allele produces the striped tabby pattern.  The recessive m allele produces the classic spiral tabby pattern. If you cross a cat that is \textit{AaMm} with a cat that is \textit{AaMm}, the expected phenotype ratio is

\begin{choices}
\choice 1 striped tabby with orange hair : 3 striped tabby with black hair : 3 classic tabby with orange hair : 1 classic tabby with black hair.
\correctchoice 9 striped tabby with orange hair : 3 striped tabby with black hair : 3 classic tabby with orange hair : 1 classic tabby with black hair. 
\choice 3 striped tabby with orange hair : 3 striped tabby with black hair : 1 classic tabby with orange hair : 1 classic tabby with black hair.
\choice 1 striped tabby with orange hair : 2 striped tabby with black hair : 2 classic tabby with orange hair : 1 classic tabby with black hair.
\choice 1 striped tabby with orange hair : 1 striped tabby with black hair : 1 classic tabby with orange hair : 1 classic tabby with black hair.
\end{choices}

\question
Based on Gregor Mendel's work, we know that traits are

\begin{choices}
\choice determined by independent assortment.
\choice determined by random segregation.
\choice inherited through the non-random passing of genes from parents to offspring.
\correctchoice inherited through the random passing of alleles from parents to offspring.
\choice determined by a random combination of dominant and recessive alleles.
\end{choices}

\newpage

\question
Assume one individual is \textit{Aa} for the first trait, and \textit{tt} for the second trait. The second individual is \textit{aa} for the first trait and \textit{Tt} for the second trait. The two traits are independent.  What is the probability of obtaining an offspring that is \textit{AaTt}? (Hint: Remember your rules for probability. You do not have to solve this using a dihybrid cross.)

\begin{choices}
\choice 0.625
\correctchoice 0.25
\choice 0.75
\choice 0.125
\choice 0.5
\end{choices}

\question
In genetics, the principles of probability can be used to

\begin{choices}
\choice state with certainty the traits of the first born child.
\choice decide which organisms are best to use in genetic crosses.
\correctchoice predict the traits of the offspring produced by genetic crosses.
\choice predict the traits of the parents used in genetic crosses.
\choice determine the actual outcomes of genetic crosses.
\end{choices}

\question
A crime suspect has alleles 9 and 12 for the CSF1PO short tandem repeat. What is the probability that someone else at random would have these same two alleles?  

\begin{multicols}{2}
\begin{choices}
\choice $1/120$.
\correctchoice $1/60$.
\choice $1/403$.
\choice $1/40$.
\choice $1/3$.
\choice $1/43$.
\end{choices}
\columnbreak
\begin{tabular}{lr}
%\begin{tabular{lr@{/}l)
\toprule
allele & frequency\\
\midrule
9 & $1/40$\\
10 & $1/5$\\
11 & $1/3$\\
12 & $1/3$\\
13 & $1/10$\\
\bottomrule
\end{tabular}
\end{multicols}

\question
DNA strands are

\begin{choices}
\correctchoice antiparallel. $5'$ matches with $3'$.
\choice antiparallel. $5'$ matches with $5'$.
\choice parallel. $3'$ matches with $3'$.
\choice complementary. $5'$ matches with $3'$.
\choice complementary. $5'$ matches with $5'$.
\choice parallel. $5'$ matches with $5'$.
\end{choices}

\question
A type of mutation that occurs when a piece of parasitic DNA inserts a copy of itself on another chromosome is called:

\begin{choices}
\choice a point mutation
\choice a frameshift mutation
\choice gene duplication
\choice an intron
\correctchoice a jumping gene
\end{choices}

\question
Danny Ramos Gomez, his brother Larry, and his sister Jaime, all have hypertrichosis. Their disease is a sex-linked, dominant trait with alleles \textit{H} and \textit{h}. Based on his family pedigree, what is the genotype of Danny's mother?
 
\begin{multicols}{2}
\begin{choices}
\choice All three genotypes are possible for Danny's mother.
\correctchoice X$^H$X$^H$ or X$^H$X$^h$
\choice X$^h$X$^h$
\choice X$^H$X$^h$
\choice X$^H$X$^H$
\end{choices}
\columnbreak
\includegraphics[width=0.4\textwidth]{gomez_pedigree}
\end{multicols}

\question
Duchene's muscular dystrophy (DMD) is a sex-linked recessive trait. Neither Betty nor her husband Bob has the trait. Their son, Buddy, was born with DMD. What is the chance that their next son will have DMD?

\begin{choices}
\choice 25\%.
\choice 75\%.
\choice 100\%.
\correctchoice 50\%.
\choice 0\%.
\end{choices}
